
\fchapter{Ej 16 Pág 55}
\fsection[\textcolor{waveBlue2}{INDICA}]{\boldit{
Indica si las siguientes acciones llevadas a cabo por empresas contribuyen o no a la erradicación de la pobreza. Justifica tu respuesta.
}} 
\begin{solution}

	\begin{enumerate}[label=\alph*)]
		\item{\boldit{Ofrecer a sus trabajadores un sueldo digno.}}
			\vspace{0,5em}

		Sí, contribuye. Un sueldo digno permite a los trabajadores cubrir sus necesidades básicas, como alimentación, vivienda y 
		educación, lo que reduce la pobreza individual y familiar. Además, mejora la productividad y el bienestar general, 
		rompiendo ciclos de pobreza intergeneracional.

			\vspace{0,5em}
		\item{\boldit{Conceder becas a estudiantes de familias con bajos ingresos.}}
			\vspace{0,5em}
			
		Sí, contribuye. Las becas facilitan el acceso a la educación superior o técnica, lo que aumenta las oportunidades 
		laborales y de ingresos futuros para personas de entornos vulnerables. Esto ayuda a erradicar la pobreza a largo 
		plazo al promover la movilidad social.
			\vspace{0,5em}
		\item{\boldit{Establecer sus fábricas en los países más pobres.}}
			\vspace{0,5em}

		Depende, pero generalmente sí contribuye (con condiciones). Si la empresa genera empleo local, transfiere tecnología y 
		respeta derechos laborales, puede impulsar el desarrollo económico y reducir la pobreza. Sin embargo, si explota mano 
		de obra barata o causa daños ambientales sin beneficios reales para la comunidad, no contribuye e incluso podría agravarla. 
		El impacto positivo requiere prácticas éticas y sostenibles.

			\vspace{0,5em}
		\item{\boldit{Conceder ayudas y microcréditos a personas de colectivos vulnerables.}}
			\vspace{0,5em}

			Sí, contribuye. Los microcréditos y ayudas permiten a individuos emprender negocios o cubrir emergencias,
			fomentando la autonomía económica. Programas como estos han demostrado éxito en modelos como el de Grameen Bank,
			reduciendo la pobreza al empoderar a mujeres y minorías.
			\vspace{0,5em}
		\item{\boldit{Establecer relaciones comerciales con pequeñas y medianas empresas.}}
			\vspace{0,5em}

			Sí, contribuye. Al comprar o colaborar con PYMES, las grandes empresas integran a proveedores locales en cadenas de
			valor, generando empleo y crecimiento económico en comunidades pobres. Esto fortalece economías locales y reduce 
			desigualdades, como en modelos de comercio justo.
	
	\end{enumerate}

\end{solution}
